% ============================================================
%  MANUEL D'UTILISATION -- DocuSense
%  Projet PFE Licence LMI -- IFAG 2025-2026
%  Auteur : Boumehdi Rym
% ============================================================
% Compiler avec : xelatex manuel_utilisation_docusense.tex

\documentclass[12pt, a4paper]{article}

% ─── Moteur XeLaTeX ───────────────────────────────────────
\usepackage{fontspec}
\usepackage{polyglossia}
\setmainlanguage{french}
\setmainfont{Times New Roman}
\setsansfont{Arial}
\setmonofont[Scale=0.85]{Courier New}

% ─── Mise en page ─────────────────────────────────────────
\usepackage[top=2.5cm, bottom=2.5cm, left=2.8cm, right=2.8cm, headheight=15pt]{geometry}
\usepackage{setspace}
\onehalfspacing
\setlength{\parindent}{0cm}
\setlength{\parskip}{6pt}

% ─── En-têtes / pieds de page ─────────────────────────────
\usepackage{fancyhdr}
\pagestyle{fancy}
\fancyhf{}
\fancyhead[L]{\small\textit{Manuel d'utilisation -- DocuSense}}
\fancyhead[R]{\small\textit{IFAG 2025-2026}}
\fancyfoot[C]{\thepage}
\renewcommand{\headrulewidth}{0.4pt}

% ─── Couleurs ─────────────────────────────────────────────
\usepackage[dvipsnames,svgnames,x11names]{xcolor}
\definecolor{DSblue}{HTML}{185FA5}
\definecolor{DSdark}{HTML}{0C2340}
\definecolor{DSlight}{HTML}{EFF6FF}
\definecolor{DSgreen}{HTML}{27AE60}
\definecolor{DSorange}{HTML}{E67E22}
\definecolor{DSgray}{HTML}{7F8C8D}
\definecolor{DSlightgray}{HTML}{F2F3F4}
\definecolor{DSpurple}{HTML}{8E44AD}

% ─── Graphiques ───────────────────────────────────────────
\usepackage{graphicx}
\usepackage{float}
\usepackage{caption}
\captionsetup{font=small, labelfont=bf, labelsep=period}

% ─── Boîtes ───────────────────────────────────────────────
\usepackage{tcolorbox}
\tcbuselibrary{skins, breakable}

\newtcolorbox{astuce}[1][Astuce]{
  colback=DSlight, colframe=DSblue, fonttitle=\bfseries\small,
  title=\faIcon{lightbulb}\quad #1, breakable, enhanced,
  left=6pt, right=6pt, top=4pt, bottom=4pt}

\newtcolorbox{attention}[1][Attention]{
  colback=orange!8, colframe=DSorange, fonttitle=\bfseries\small,
  title=#1, breakable, enhanced,
  left=6pt, right=6pt, top=4pt, bottom=4pt}

\newtcolorbox{etape}[2][]{
  colback=DSlightgray, colframe=DSblue!60, fonttitle=\bfseries,
  title={\textcolor{DSblue}{\Large\bfseries#2}},
  breakable, enhanced, left=8pt, right=8pt, top=6pt, bottom=6pt,
  attach boxed title to top left={yshift=-2mm, xshift=8mm},
  boxed title style={colback=DSblue, colframe=DSblue}}

% ─── Placeholder capture d'écran ──────────────────────────
\usepackage{tikz}
\newcommand{\capture}[2][10cm]{%
  \begin{center}
    \begin{tikzpicture}
      \draw[DSgray, dashed, rounded corners=4pt, line width=1pt]
        (0,0) rectangle (#1, -5.5cm);
      \node[DSgray] at (#1/2, -2.75cm) {%
        \begin{minipage}{0.85\linewidth}\centering
          \textit{\small [Capture d'écran : #2]}
        \end{minipage}};
    \end{tikzpicture}
  \end{center}}

% ─── Listes ───────────────────────────────────────────────
\usepackage{enumitem}
\setlist[itemize]{leftmargin=1.5em, itemsep=2pt}
\setlist[enumerate]{leftmargin=1.8em, itemsep=2pt}

% ─── Hyperliens ───────────────────────────────────────────
\usepackage[unicode=true, colorlinks=true, linkcolor=DSblue,
  urlcolor=DSblue, bookmarks=true,
  pdftitle={Manuel d'utilisation -- DocuSense},
  pdfauthor={Boumehdi Rym}]{hyperref}

% ─── Divers ───────────────────────────────────────────────
\usepackage{booktabs}
\usepackage{array}
\usepackage{tabularx}
\usepackage{multirow}
\usepackage{amsmath}
\usepackage{pifont}
\newcommand{\cmark}{\textcolor{DSgreen}{\ding{51}}}
\newcommand{\xmark}{\textcolor{red}{\ding{55}}}

% ─── Numérotation sections ─────────────────────────────────
\setcounter{secnumdepth}{3}
\setcounter{tocdepth}{2}

% ──────────────────────────────────────────────────────────
\begin{document}

% ═══════════════════════════════════════════════════════════
%  PAGE DE TITRE
% ═══════════════════════════════════════════════════════════
\begin{titlepage}
\thispagestyle{empty}
\begin{tikzpicture}[remember picture, overlay]
  % Bandeau bleu en haut
  \fill[DSdark] (current page.north west) rectangle
    ([yshift=-3.5cm]current page.north east);
  % Bande bleue claire
  \fill[DSblue] ([yshift=-3.5cm]current page.north west) rectangle
    ([yshift=-4cm]current page.north east);
  % Pied de page
  \fill[DSlightgray] (current page.south west) rectangle
    ([yshift=2cm]current page.south east);
\end{tikzpicture}

\vspace*{0.5cm}
\begin{center}
  {\color{white}\sffamily\Large\bfseries IFAG -- Institut de Formation en
  Algérie de Gestion}\\[0.3cm]
  {\color{white}\sffamily\large Licence LMI -- Année universitaire 2025-2026}
\end{center}

\vspace{3cm}

\begin{center}
  % Logo / nom plateforme
  \begin{tikzpicture}
    \node[rounded corners=8pt, fill=DSblue, text=white,
          font=\sffamily\Huge\bfseries,
          inner xsep=20pt, inner ysep=12pt] {DocuSense};
  \end{tikzpicture}

  \vspace{0.8cm}
  {\sffamily\huge\bfseries\color{DSdark} Manuel d'Utilisation}\\[0.4cm]
  {\sffamily\large\color{DSgray} Guide pas à pas avec scénario de démonstration}

  \vspace{1.5cm}
  \begin{tcolorbox}[colback=DSlight, colframe=DSblue, width=0.75\linewidth,
    halign=center, boxrule=1pt]
    \textbf{Plateforme intelligente de génération de\\
    documentation technique gamifiée}\\[4pt]
    {\small Propulsée par l'IA générative -- Google Gemini 2.0 Flash}
  \end{tcolorbox}

  \vspace{2cm}
  {\large\bfseries Réalisé par :} {\large Boumehdi Rym}\\[0.3cm]
  {\large\bfseries Encadrant :} {\large Mr. Abbas}\\[0.6cm]
  {\large\color{DSgray} Juin 2026}
\end{center}
\end{titlepage}

% ═══════════════════════════════════════════════════════════
%  TABLE DES MATIÈRES
% ═══════════════════════════════════════════════════════════
\newpage
\tableofcontents
\newpage

% ═══════════════════════════════════════════════════════════
%  INTRODUCTION
% ═══════════════════════════════════════════════════════════
\section{Introduction}

Ce manuel décrit, étape par étape, comment utiliser la plateforme
\textbf{DocuSense}. Il s'adresse à deux types d'utilisateurs :

\begin{itemize}
  \item \textbf{Le Développeur} : crée des manuels techniques,
        génère des étapes automatiquement via l'IA, publie le contenu
        et crée des quiz.
  \item \textbf{L'Utilisateur final (End User)} : consulte les manuels
        publiés, suit les tutoriels étape par étape, répond aux quiz
        et remporte des badges.
\end{itemize}

\begin{tcolorbox}[colback=DSlight, colframe=DSblue, title=\textbf{Accès à la plateforme}]
  Ouvrez votre navigateur web (Google Chrome, Firefox, Edge) et
  accédez à l'adresse suivante :\\[4pt]
  \texttt{\textbf{http://localhost:3000}} \quad (environnement local)
\end{tcolorbox}

\vspace{0.3cm}

% ═══════════════════════════════════════════════════════════
%  SECTION 1 : PAGE D'ACCUEIL
% ═══════════════════════════════════════════════════════════
\section{Page d'accueil}

La page d'accueil est la première interface visible par tout visiteur.
Elle présente la plateforme DocuSense, ses fonctionnalités principales
et les boutons d'accès à l'inscription ou à la connexion.

\capture[14cm]{Page d'accueil -- vue générale de DocuSense}

\vspace{0.3cm}

La page d'accueil contient les éléments suivants :

\begin{itemize}
  \item \textbf{Barre de navigation} : liens vers les rubriques
        \textit{Fonctionnalités}, \textit{Tarifs}, \textit{Documentation}
        et \textit{À propos}.
  \item \textbf{Section héro} : titre principal, description et deux
        boutons d'action (\textit{Commencer gratuitement} /
        \textit{Voir la démo}).
  \item \textbf{Section fonctionnalités} : présentation des trois
        piliers de la plateforme (IA générative, Gamification,
        Ultra rapide).
  \item \textbf{Section tarifs} : offre 100\,\% gratuite.
  \item \textbf{Section À propos} : contexte du projet IFAG / LMI.
\end{itemize}

\capture[14cm]{Section "Fonctionnalités" de la page d'accueil}

% ═══════════════════════════════════════════════════════════
%  SECTION 2 : INSCRIPTION
% ═══════════════════════════════════════════════════════════
\section{Inscription (Créer un compte)}

\subsection{Accéder à la page d'inscription}

Depuis la page d'accueil, cliquez sur le bouton
\textbf{« Commencer gratuitement »} ou sur le lien \textbf{« S'inscrire »}
dans la barre de navigation.

\capture[14cm]{Bouton "Commencer gratuitement" sur la page d'accueil}

\subsection{Remplir le formulaire d'inscription}

La page d'inscription affiche un formulaire composé des champs suivants :

\begin{enumerate}
  \item \textbf{Nom complet} : saisissez votre prénom et nom.
  \item \textbf{Adresse e-mail} : entrez une adresse e-mail valide.
  \item \textbf{Mot de passe} : choisissez un mot de passe sécurisé
        (minimum 6 caractères).
  \item \textbf{Rôle} : sélectionnez votre profil :
    \begin{itemize}
      \item \textit{Développeur} -- pour créer et gérer des manuels.
      \item \textit{Utilisateur} -- pour consulter les manuels publiés.
    \end{itemize}
\end{enumerate}

\capture[14cm]{Formulaire d'inscription avec sélection du rôle}

\subsection{Valider l'inscription}

Après avoir rempli tous les champs, cliquez sur le bouton
\textbf{« S'inscrire »}. Un message de confirmation s'affiche et vous
êtes redirigé(e) vers la page de connexion.

\begin{tcolorbox}[colback=orange!8, colframe=DSorange,
  title=\textbf{Remarque}]
  Vous pouvez également vous inscrire rapidement via votre compte
  \textbf{Google} en cliquant sur le bouton \textit{« Continuer avec Google »}.
\end{tcolorbox}

% ═══════════════════════════════════════════════════════════
%  SECTION 3 : CONNEXION
% ═══════════════════════════════════════════════════════════
\section{Connexion (Se connecter)}

\subsection{Accéder à la page de connexion}

Depuis la page d'accueil, cliquez sur le bouton \textbf{« Se connecter »}
dans la barre de navigation.

\subsection{Saisir les identifiants}

\begin{enumerate}
  \item Entrez votre \textbf{adresse e-mail} et votre \textbf{mot de passe}.
  \item Cliquez sur \textbf{« Se connecter »}.
\end{enumerate}

\capture[14cm]{Page de connexion avec formulaire email et mot de passe}

Selon le rôle associé à votre compte, vous serez automatiquement
redirigé(e) vers :
\begin{itemize}
  \item \textbf{Le tableau de bord (Dashboard)} si vous êtes Développeur.
  \item \textbf{L'espace utilisateur (End User)} si vous êtes Utilisateur.
\end{itemize}

% ═══════════════════════════════════════════════════════════
%  SECTION 4 : ESPACE DÉVELOPPEUR
% ═══════════════════════════════════════════════════════════
\section{Espace Développeur -- Tableau de bord}

L'espace développeur est le cœur de la création de contenu sur DocuSense.
Il permet de créer, modifier, publier et archiver des manuels techniques.

\capture[14cm]{Tableau de bord -- vue générale avec liste des manuels}

\subsection{Créer un nouveau manuel}

\textbf{Étape 1 -- Cliquer sur « Nouveau manuel »}

Dans la barre latérale ou en haut de la liste, cliquez sur le bouton
\textbf{« + Nouveau manuel »}.

\capture[14cm]{Bouton "Nouveau manuel" dans le tableau de bord}

\textbf{Étape 2 -- Renseigner les informations du manuel}

Un formulaire apparaît avec les champs suivants :
\begin{itemize}
  \item \textbf{Titre} : nom du manuel (ex. : \textit{Guide d'installation
        de Node.js}).
  \item \textbf{Description} : brève description du sujet traité.
\end{itemize}

Cliquez sur \textbf{« Créer »} pour valider.

\capture[14cm]{Formulaire de création d'un nouveau manuel}

\subsection{Générer les étapes automatiquement avec l'IA}

DocuSense propose plusieurs méthodes de génération automatique d'étapes
grâce à l'intelligence artificielle (Google Gemini 2.0 Flash).

\subsubsection{Méthode 1 : Génération par description texte}

\begin{enumerate}
  \item Ouvrez le manuel créé.
  \item Dans la section \textbf{« Générer avec l'IA »}, saisissez une
        description de votre processus (ex. : \textit{«~Comment installer
        Node.js sur Windows~»}).
  \item Cliquez sur \textbf{« Suggérer les étapes »}.
  \item L'IA génère automatiquement 5 étapes structurées.
\end{enumerate}

\capture[14cm]{Interface de génération IA par description texte}

\capture[14cm]{Résultat : 5 étapes générées automatiquement par l'IA}

\subsubsection{Méthode 2 : Analyse d'une capture d'écran}

\begin{enumerate}
  \item Cliquez sur l'onglet \textbf{« Analyser une image »}.
  \item Chargez une capture d'écran (format PNG ou JPG).
  \item Cliquez sur \textbf{« Analyser »}.
  \item L'IA analyse l'image et génère les étapes correspondantes.
\end{enumerate}

\capture[14cm]{Upload d'une capture d'écran pour analyse IA}

\subsubsection{Méthode 3 : Analyse d'une URL de site web}

\begin{enumerate}
  \item Cliquez sur l'onglet \textbf{« Analyser une URL »}.
  \item Collez l'URL du site web à documenter.
  \item Cliquez sur \textbf{« Analyser »}.
  \item L'IA génère les étapes de navigation sur ce site.
\end{enumerate}

\capture[14cm]{Interface d'analyse par URL}

\subsubsection{Méthode 4 : Analyse d'une vidéo}

\begin{enumerate}
  \item Cliquez sur l'onglet \textbf{« Analyser une vidéo »}.
  \item Chargez une vidéo de démonstration (MP4).
  \item Cliquez sur \textbf{« Analyser »}.
  \item L'IA extrait les étapes à partir du contenu vidéo.
\end{enumerate}

\capture[14cm]{Upload d'une vidéo pour analyse IA}

\subsection{Modifier les étapes manuellement}

Après la génération automatique (ou manuellement), vous pouvez modifier
chaque étape :
\begin{enumerate}
  \item Cliquez sur une étape pour l'éditer.
  \item Modifiez le \textbf{titre} et la \textbf{description}.
  \item Ajoutez ou remplacez le \textbf{média} associé (image/vidéo).
  \item Cliquez sur \textbf{« Enregistrer »}.
\end{enumerate}

\capture[14cm]{Éditeur d'étape -- modification du titre et description}

\subsection{Créer un quiz pour une étape}

Chaque étape peut comporter un quiz de validation des connaissances.

\begin{enumerate}
  \item Sélectionnez une étape dans le manuel.
  \item Cliquez sur \textbf{« Ajouter un quiz »}.
  \item Remplissez les champs :
    \begin{itemize}
      \item \textbf{Question} : posez une question sur le contenu de l'étape.
      \item \textbf{Options} : saisissez 4 propositions de réponse (A, B, C, D).
      \item \textbf{Bonne réponse} : indiquez la réponse correcte.
      \item \textbf{Points} : attribuez un nombre de points (ex. : 10 pts).
    \end{itemize}
  \item Cliquez sur \textbf{« Créer le quiz »}.
\end{enumerate}

\capture[14cm]{Formulaire de création d'un quiz pour une étape}

\subsection{Publier un manuel}

Une fois le manuel complet, publiez-le pour le rendre accessible
aux utilisateurs finaux.

\begin{enumerate}
  \item Dans la liste des manuels, localisez le manuel à publier.
  \item Cliquez sur le bouton \textbf{« Publier »} (icône de publication).
  \item Le statut du manuel passe à \textbf{« Publié »}.
\end{enumerate}

\capture[14cm]{Manuel avec statut "Publié" visible dans le tableau de bord}

\begin{tcolorbox}[colback=DSlight, colframe=DSblue, title=\textbf{Information}]
  Un manuel publié est immédiatement visible dans l'espace utilisateur
  (End User). Vous pouvez le dépublier à tout moment en cliquant à nouveau
  sur le bouton de publication.
\end{tcolorbox}

\subsection{Exporter un manuel en PDF}

\begin{enumerate}
  \item Ouvrez le manuel souhaité.
  \item Cliquez sur le bouton \textbf{« Exporter PDF »}.
  \item Le fichier PDF est automatiquement téléchargé sur votre ordinateur.
\end{enumerate}

\capture[14cm]{Bouton "Exporter PDF" et aperçu du fichier généré}

\subsection{Archiver et restaurer un manuel}

\begin{itemize}
  \item \textbf{Archiver} : cliquez sur l'icône d'archivage à côté du
        manuel. Il disparaît de la liste principale et passe dans
        les \textit{Archives}.
  \item \textbf{Restaurer} : accédez à la section \textit{Archives},
        cliquez sur \textbf{« Restaurer »} pour remettre le manuel
        dans la liste active.
\end{itemize}

\capture[14cm]{Section Archives avec bouton "Restaurer"}

% ═══════════════════════════════════════════════════════════
%  SECTION 5 : ESPACE UTILISATEUR FINAL
% ═══════════════════════════════════════════════════════════
\section{Espace Utilisateur Final (End User)}

L'espace utilisateur final permet de consulter les manuels publiés,
de suivre les tutoriels, de répondre aux quiz et de gagner des récompenses.

\capture[14cm]{Espace utilisateur -- liste des manuels disponibles}

\subsection{Parcourir les manuels disponibles}

\begin{enumerate}
  \item Connectez-vous avec un compte de type \textbf{Utilisateur}.
  \item La page principale affiche tous les manuels publiés.
  \item Cliquez sur un manuel pour commencer à le consulter.
\end{enumerate}

\capture[14cm]{Liste des manuels publiés avec titre et description}

\subsection{Suivre un tutoriel étape par étape}

\begin{enumerate}
  \item Après avoir ouvert un manuel, les étapes s'affichent dans l'ordre.
  \item Lisez la description de chaque étape.
  \item Consultez le média associé (image ou vidéo) si disponible.
  \item Cliquez sur \textbf{« Étape suivante »} pour progresser.
\end{enumerate}

\capture[14cm]{Vue d'une étape avec description et média associé}

\capture[14cm]{Navigation entre les étapes (boutons Précédent / Suivant)}

\subsection{Répondre aux quiz}

Lorsqu'une étape comporte un quiz, il s'affiche automatiquement
après avoir lu le contenu.

\begin{enumerate}
  \item Lisez la \textbf{question} affichée.
  \item Sélectionnez une des quatre options proposées (A, B, C ou D).
  \item Cliquez sur \textbf{« Valider ma réponse »}.
  \item Le résultat s'affiche :
    \begin{itemize}
      \item \cmark\ \textbf{Bonne réponse} : des points sont ajoutés
            à votre score.
      \item \xmark\ \textbf{Mauvaise réponse} : la bonne réponse
            vous est indiquée.
    \end{itemize}
\end{enumerate}

\capture[14cm]{Interface du quiz avec les 4 options de réponse}

\capture[14cm]{Résultat après validation -- bonne ou mauvaise réponse}

\subsection{Badges et système de récompenses}

DocuSense intègre un système de \textbf{gamification} basé sur
l'accumulation de points :

\begin{center}
\begin{tabular}{clcc}
\toprule
\textbf{Badge} & \textbf{Nom} & \textbf{Points requis} & \textbf{Couleur} \\
\midrule
\ding{72} & Bronze & 50 pts  & \textcolor{DSorange}{\rule{1cm}{6pt}} \\
\ding{72} & Silver & 100 pts & \textcolor{DSgray}{\rule{1cm}{6pt}} \\
\ding{72} & Gold   & 200 pts & \textcolor{yellow!60!DSorange}{\rule{1cm}{6pt}} \\
\bottomrule
\end{tabular}
\end{center}

\vspace{0.3cm}
Les badges obtenus s'affichent dans votre profil utilisateur.

\capture[14cm]{Profil utilisateur affichant les badges obtenus}

\subsection{Classement (Leaderboard)}

Le \textbf{leaderboard} affiche le classement des 10 meilleurs
utilisateurs selon leur score total.

\begin{enumerate}
  \item Accédez à la section \textbf{« Classement »} dans le menu.
  \item Consultez votre position dans le tableau.
\end{enumerate}

\capture[14cm]{Tableau du classement -- Top 10 des utilisateurs}

% ═══════════════════════════════════════════════════════════
%  SECTION 6 : SCÉNARIO DE DÉMONSTRATION
% ═══════════════════════════════════════════════════════════
\newpage
\section{Scénario de Démonstration}

\begin{tcolorbox}[colback=DSlight, colframe=DSblue,
  title=\textbf{Description du scénario}]
  Ce scénario illustre le cycle complet d'utilisation de DocuSense :
  un développeur crée un manuel technique via l'IA, puis un utilisateur
  final le consulte, répond aux quiz et remporte un badge.
\end{tcolorbox}

\vspace{0.5cm}

\subsection*{Acteurs}
\begin{itemize}
  \item \textbf{Rym} -- Développeur, créatrice du manuel.
  \item \textbf{Karim} -- Utilisateur final, apprenant.
\end{itemize}

\vspace{0.3cm}

%----- Étape A -----
\begin{tcolorbox}[colback=DSlightgray, colframe=DSblue, breakable,
  title={\textbf{Étape A -- Rym s'inscrit en tant que Développeur}}]

Rym ouvre DocuSense dans son navigateur et clique sur
\textit{« Commencer gratuitement »}. Elle remplit le formulaire
d'inscription avec son nom, son e-mail et choisit le rôle
\textbf{Développeur}. Elle valide et se retrouve sur le tableau de bord.

\end{tcolorbox}
\capture[14cm]{Scénario -- Rym complète le formulaire d'inscription}

\vspace{0.3cm}

%----- Étape B -----
\begin{tcolorbox}[colback=DSlightgray, colframe=DSblue, breakable,
  title={\textbf{Étape B -- Rym crée un manuel "Comment utiliser DocuSense"}}]

Depuis son tableau de bord, Rym clique sur \textit{« + Nouveau manuel »}.
Elle saisit le titre \textit{« Comment utiliser DocuSense »} et une courte
description. Elle valide la création.

\end{tcolorbox}
\capture[14cm]{Scénario -- Rym crée le manuel "Comment utiliser DocuSense"}

\vspace{0.3cm}

%----- Étape C -----
\begin{tcolorbox}[colback=DSlightgray, colframe=DSblue, breakable,
  title={\textbf{Étape C -- Rym génère les étapes par IA (description texte)}}]

Dans le manuel créé, Rym tape dans le champ IA :
\textit{«~Guide pas à pas pour utiliser la plateforme DocuSense~»}.
Elle clique sur \textit{« Suggérer les étapes »}. En quelques secondes,
Gemini génère \textbf{5 étapes structurées} : inscription, connexion,
création de manuel, publication et quiz.

\end{tcolorbox}
\capture[14cm]{Scénario -- 5 étapes générées par Gemini IA en quelques secondes}

\vspace{0.3cm}

%----- Étape D -----
\begin{tcolorbox}[colback=DSlightgray, colframe=DSblue, breakable,
  title={\textbf{Étape D -- Rym ajoute un quiz sur l'étape 1}}]

Rym sélectionne la première étape \textit{«~Inscription~»} et clique
sur \textit{« Ajouter un quiz »}. Elle crée la question :
\textit{«~Quel rôle choisir pour créer des manuels ?~»} avec les options
A) Développeur, B) Utilisateur, C) Administrateur, D) Visiteur.
La bonne réponse est \textbf{A -- Développeur}, valant 10 points.

\end{tcolorbox}
\capture[14cm]{Scénario -- Rym crée un quiz sur l'étape 1}

\vspace{0.3cm}

%----- Étape E -----
\begin{tcolorbox}[colback=DSlightgray, colframe=DSblue, breakable,
  title={\textbf{Étape E -- Rym publie le manuel}}]

Une fois toutes les étapes et quiz configurés, Rym clique sur
\textit{« Publier »}. Le statut du manuel passe à \textbf{Publié}
et il devient accessible aux utilisateurs finaux.

\end{tcolorbox}
\capture[14cm]{Scénario -- Rym publie le manuel, statut "Publié" affiché}

\vspace{0.3cm}

%----- Étape F -----
\begin{tcolorbox}[colback=DSlightgray, colframe=DSgreen!60, breakable,
  title={\textbf{Étape F -- Karim s'inscrit en tant qu'Utilisateur}}]

Karim ouvre DocuSense, s'inscrit avec le rôle \textbf{Utilisateur}
et accède à l'espace End User. Il voit immédiatement le manuel
\textit{« Comment utiliser DocuSense »} publié par Rym.

\end{tcolorbox}
\capture[14cm]{Scénario -- Karim voit le manuel de Rym dans l'espace utilisateur}

\vspace{0.3cm}

%----- Étape G -----
\begin{tcolorbox}[colback=DSlightgray, colframe=DSgreen!60, breakable,
  title={\textbf{Étape G -- Karim consulte le tutoriel étape par étape}}]

Karim clique sur le manuel et parcourt les 5 étapes une par une.
Il lit les descriptions et regarde les médias associés à chaque étape.
Il navigue avec les boutons \textit{Précédent} et \textit{Suivant}.

\end{tcolorbox}
\capture[14cm]{Scénario -- Karim navigue entre les étapes du tutoriel}

\vspace{0.3cm}

%----- Étape H -----
\begin{tcolorbox}[colback=DSlightgray, colframe=DSgreen!60, breakable,
  title={\textbf{Étape H -- Karim répond au quiz et gagne des points}}]

À la fin de l'étape 1, le quiz apparaît. Karim sélectionne
\textbf{A -- Développeur} et clique sur \textit{« Valider »}.
La bonne réponse est confirmée et \textbf{+10 points} sont ajoutés
à son score. Il complète les autres quiz et accumule au total 60 points.

\end{tcolorbox}
\capture[14cm]{Scénario -- Karim valide sa réponse et gagne +10 points}

\vspace{0.3cm}

%----- Étape I -----
\begin{tcolorbox}[colback=DSlightgray, colframe=DSorange!80, breakable,
  title={\textbf{Étape I -- Karim remporte le badge Bronze}}]

Avec 60 points accumulés (seuil Bronze = 50 pts), le badge
\textbf{Bronze} s'affiche automatiquement sur le profil de Karim.
Il consulte le \textbf{leaderboard} et voit son nom figurer dans
le classement.

\end{tcolorbox}
\capture[14cm]{Scénario -- Badge Bronze attribué à Karim sur son profil}

\capture[14cm]{Scénario -- Karim apparaît dans le leaderboard Top 10}

% ═══════════════════════════════════════════════════════════
%  SECTION 7 : RÉCAPITULATIF DES FONCTIONNALITÉS
% ═══════════════════════════════════════════════════════════
\newpage
\section{Récapitulatif des fonctionnalités}

\begin{center}
\renewcommand{\arraystretch}{1.4}
\begin{tabularx}{\linewidth}{lXc}
\toprule
\textbf{Fonctionnalité} & \textbf{Description} & \textbf{Profil} \\
\midrule
Inscription / Connexion & Créer un compte et se connecter (email ou Google OAuth) & Tous \\
\midrule
Création de manuel & Créer un manuel avec titre et description & Développeur \\
Génération IA (texte) & Générer des étapes via une description textuelle & Développeur \\
Génération IA (image) & Analyser une capture d'écran pour générer des étapes & Développeur \\
Génération IA (URL) & Analyser un site web pour générer des étapes & Développeur \\
Génération IA (vidéo) & Analyser une vidéo pour générer des étapes & Développeur \\
Édition des étapes & Modifier manuellement le titre, description et médias & Développeur \\
Création de quiz & Ajouter un quiz (QCM 4 options) à chaque étape & Développeur \\
Publication & Rendre un manuel accessible aux utilisateurs finaux & Développeur \\
Export PDF & Télécharger le manuel en fichier PDF & Développeur \\
Archivage & Archiver / restaurer un manuel & Développeur \\
\midrule
Consultation manuals & Parcourir les manuels publiés & Utilisateur \\
Navigation étapes & Suivre un tutoriel étape par étape & Utilisateur \\
Quiz interactif & Répondre aux quiz et gagner des points & Utilisateur \\
Badges & Obtenir des badges (Bronze / Silver / Gold) selon les points & Utilisateur \\
Leaderboard & Consulter le classement des 10 meilleurs utilisateurs & Utilisateur \\
\bottomrule
\end{tabularx}
\end{center}

% ═══════════════════════════════════════════════════════════
%  CONCLUSION
% ═══════════════════════════════════════════════════════════
\section{Conclusion}

Ce manuel d'utilisation a présenté, étape par étape, l'ensemble des
fonctionnalités de la plateforme \textbf{DocuSense}. Le scénario de
démonstration a illustré le cycle complet : de la création d'un manuel
par un développeur jusqu'à l'obtention d'un badge par un utilisateur final.

DocuSense simplifie la documentation technique en automatisant la
génération d'étapes grâce à l'IA générative, tout en engageant les
apprenants via la gamification (quiz, badges, classement).

\vspace{1cm}
\begin{tcolorbox}[colback=DSlightgray, colframe=DSgray]
  \textbf{Contact :} \texttt{boumehdirym0@gmail.com} \quad
  \textbf{Encadrant :} Mr. Abbas \quad
  \textbf{IFAG -- Licence LMI -- 2025-2026}
\end{tcolorbox}

\end{document}
